%-------------------------
% Resume in Latex
% Author : Jake Gutierrez
% Based off of: https://github.com/sb2nov/resume
% License : MIT
%------------------------

\documentclass[letterpaper,10pt]{article}

\usepackage{latexsym}
\usepackage{graphicx}
\usepackage{dirtytalk}
\usepackage{ragged2e}
\usepackage[empty]{fullpage}
\usepackage{titlesec}
\usepackage{marvosym}
\usepackage[usenames,dvipsnames]{xcolor}
\usepackage{verbatim}
\usepackage{enumitem}
\usepackage[hidelinks]{hyperref}
\usepackage{fancyhdr}
\usepackage[english]{babel}
\usepackage{tabularx}
\usepackage{tikz}
\usepackage[absolute,overlay]{textpos}
\setlength{\TPHorizModule}{1in}
\setlength{\TPVertModule}{1in}
\input{glyphtounicode}


%----------FONT OPTIONS----------
% sans-serif
% \usepackage[sfdefault]{FiraSans}
% \usepackage[sfdefault]{roboto}
% \usepackage[sfdefault]{noto-sans}
% \usepackage[default]{sourcesanspro}

% serif
% \usepackage{CormorantGaramond}
% \usepackage{charter}


\pagestyle{fancy}
\fancyhf{} % clear all header and footer fields
\fancyfoot{}
\renewcommand{\headrulewidth}{0pt}
\renewcommand{\footrulewidth}{0pt}

% Adjust margins
\addtolength{\oddsidemargin}{-0.5in}
\addtolength{\evensidemargin}{-0.5in}
\addtolength{\textwidth}{1in}
\addtolength{\topmargin}{-.5in}
\addtolength{\textheight}{1.0in}

\urlstyle{same}

\raggedbottom
\raggedright
\setlength{\tabcolsep}{0in}

% Sections formatting
\titleformat{\section}{
  \vspace{-4pt}\scshape\raggedright\large
}{}{0em}{}[\color{black}\titlerule \vspace{-5pt}]

\titleformat{\littlesection}{
  \vspace{-4pt}\scshape\raggedright
}{}{0em}{}[\color{black}\titlerule \vspace{-5pt}]

% Ensure that generate pdf is machine readable/ATS parsable
\pdfgentounicode=1

%-------------------------
% Custom commands
\newcommand{\resumeItem}[1]{
  \item\small{
    {#1 \vspace{-2pt}}
  }
}

\def\code#1{\texttt{#1}}

\newcommand{\fancyurl}[2]{\textcolor{BlueViolet}{\href{#1}{\underline{#2}}}}
\newcommand{\fancyref}[1]{\textcolor{BlueViolet}{\nameref{#1}}}

\newcommand{\resumeSubheading}[4]{
  \vspace{-2pt}\item
    \begin{tabular*}{0.97\textwidth}[t]{l@{\extracolsep{\fill}}r}
      \textbf{#1} & #2 \\
      \textit{\small#3} & \textit{\small #4} \\
    \end{tabular*}\vspace{-3pt}
}

\newcommand{\resumeLittleSubheading}[4]{
  \vspace{-2pt}\item
    \begin{tabular*}{0.97\textwidth}[t]{l@{\extracolsep{\fill}}r}
      \footnotesize{\textbf{#1}} & \footnotesize{{#2}} \\
      \textit{\footnotesize #3} & \textit{\footnotesize #4} \\
    \end{tabular*}\vspace{-3pt}
}

\newcommand{\resumeBody}[1]{
  \vspace{2pt} \\ \small{#1}
}

\newcommand{\resumeTitleSubheading}[3]{
  \vspace{-2pt}\item
    \begin{tabular*}{0.97\textwidth}[t]{l@{\extracolsep{\fill}}r}
      \textbf{#1} & #2 \\
    \end{tabular*}
    \vspace{2pt} \\
    \small{\textbf{Title}: \textit{#3}}
    \vspace{-7pt}
}

\newcommand{\resumePaperTitle}[5]{
    \vspace{2pt}
    \item
    \small{\textbf{#1}, \textit{#2}, #3, (#4) - #5} \\
}

\newcommand{\resumeSubSubheading}[2]{
    \item
    \begin{tabular*}{0.97\textwidth}{l@{\extracolsep{\fill}}r}
      \textit{\small#1} & \textit{\small #2} \\
    \end{tabular*}\vspace{-7pt}
}

\newcommand{\resumeProjectHeading}[2]{
    \item
    \begin{tabular*}{0.97\textwidth}{l@{\extracolsep{\fill}}r}
      #1 & #2 \\
    \end{tabular*}\vspace{-3pt}
}

\newcommand{\resumeSubItem}[1]{\resumeItem{#1}\vspace{-4pt}}

\renewcommand\labelitemii{$\vcenter{\hbox{\tiny$\bullet$}}$}

\newcommand{\resumeSubHeadingListStart}{\begin{itemize}[leftmargin=0.15in, label={}]\vspace{2pt}}
\newcommand{\resumeSubHeadingPublicationsListStart}{\begin{itemize}[leftmargin=0.15in, rightmargin=0.05in, label={}]}
\newcommand{\resumeSubHeadingListEnd}{\end{itemize}}
\newcommand{\resumeItemListStart}{\begin{itemize}[rightmargin=0.1in]\vspace{-3pt}}
\newcommand{\resumeItemListEnd}{\end{itemize}\vspace{-3pt}}

\newcommand{\programmingSkills}[4]{
    \textbf{#1}
    \begin{itemize}[leftmargin=0.15in, topsep=0pt, label={}]
    \small{\item
    \textit{Proficiency}: {#2} \\
    \textit{Intermediate}: {#3} \\
    \textit{Basics}: {#4}}
    \end{itemize}
}

%-------------------------------------------
%%%%%%  RESUME STARTS HERE  %%%%%%%%%%%%%%%%%%%%%%%%%%%%


\begin{document}

%----------HEADING----------
\begin{center}
        \textbf{\huge \scshape Biagio Scarciello} \\ \vspace{4pt} \large{Embedded Software \& Systems Engineer} \\ \vspace{2pt}
        \small{C/C++ $|$ Real-Time Systems $|$ Embedded Linux $|$ Railway Safety-Critical Systems} \\ \vspace{4pt}
        \small \fancyurl{mailto:biscarci@gmail.com}{biscarci@gmail.com} $|$
        +39 393 23 07 691 $|$
        \fancyurl{https://www.linkedin.com/in/biscarci90}{LinkedIn} $|$
        \fancyurl{https://github.com/biscarci}{GitHub} \\ \vspace{2pt}
        \small{Via Sant'Antonio a Capodimonte 46, 80131 Naples, Italy}
    \end{center}


\section{Profile}

\justifying
\small{Senior Embedded Software \& Systems Engineer with 8+ years of experience delivering
safety-critical railway systems across high-speed, metro, and regional platforms. Expert in
Embedded C/C++, real-time system design, and onboard communication protocols (TRDP, IPTCom,
Ethernet). Proven technical referent over the full engineering lifecycle: system architecture
definition, integration, verification \& validation, and on-site commissioning. Experienced in
embedded Linux environments and in cross-functional, international engineering teams.
Open to relocation and international assignments.}

%-----------EXPERIENCE-----------
\section{Work Experience}
  \resumeSubHeadingListStart
    \resumeSubheading
        {Software Lead --- Diagnostic Systems}{2024 -- Present}
        {Hitachi Rail STS S.p.A.}{Naples, Italy}
        \resumeItemListStart
          \resumeItem{
            \textbf{Software lead} for the diagnostic systems of the \textbf{Rome Metro (lines A/B)}
            and \textbf{Florence Tram} projects, owning the software from development through
            integration and on-site commissioning.
          }
          \resumeItem{
            \textbf{Software development support} on the \textbf{Milan Metro (lines 1, 2, 3)}
            diagnostic system, contributing to development, integration, and issue resolution.
          }
          \resumeItem{
            Led a focused stabilization phase on the Rome Metro project ahead of official customer
            testing, \textbf{resolving 230 open issues within one month} and securing the validation
            milestones on schedule.
          }
          \resumeItem{
            Led the full software lifecycle of onboard diagnostic systems, developing the HMI and
            optimizing the embedded backend framework in Embedded C.
          }
        \resumeItemListEnd

    \resumeSubheading
        {System Engineer --- Event Recorder \& Safety-Critical Subsystems}{2022 -- 2024}
        {Hitachi Rail STS S.p.A.}{Naples, Italy}
        \resumeItemListStart
          \resumeItem{
            System owner of the \textbf{Event Recorder} product line: defined the complete system
            architecture and requirements, and drove the solution from specification through
            integration, validation, and customer acceptance for the metro networks of
            \textbf{Milan, Rome, and Washington (WMATA)}.
          }
          \resumeItem{
            System Engineer referent for Bogie Control System communication (TCMS integration) using
            message data protocol, defining the communication architecture and coordinating
            cross-functional validation activities across software, hardware, and RAMS teams.
          }
          \resumeItem{
            Technical referent for onboard safety-critical subsystems --- Event Recorder, Dead Man,
            and Zero Speed detection --- covering requirements analysis, traceability (DOORS),
            interface definition, and verification strategy.
          }
          \resumeItem{
            Coordinated integration, troubleshooting, and commissioning of diagnostic subsystems in
            operational environments, improving the reliability of diagnostic data communication
            between subsystems and TCMS.
          }
        \resumeItemListEnd

    \resumeSubheading
        {Embedded Software Developer --- Regional \& High-Speed Platforms}{Apr. 2018 -- 2022}
        {Hitachi Rail STS S.p.A.}{Naples, Italy}
        \resumeItemListStart
          \resumeItem{
            Developed and commissioned the diagnostic systems for the regional fleets
            \textbf{Trenitalia Rock}, \textbf{Trenitalia Blues}, and \textbf{Ferrovie Nord Milano},
            and contributed to the high-speed platforms (\textit{ETR1000 Frecciarossa, ETR700}).
          }
          \resumeItem{
            Developed and extended Condition-Based Maintenance (CBM) software in Embedded C, and
            designed diagnostic algorithms in C and MATLAB/Simulink for onboard systems.
          }
          \resumeItem{
            Enhanced and refactored the embedded diagnostic software framework, improving
            maintainability and reuse across train platforms.
          }
          \resumeItem{
            Implemented multi-vehicle composition management logic (backend and frontend) and
            designed diagnostic HMI components in HTML/JavaScript.
          }
          \resumeItem{
            Improved the reliability and stability of distributed onboard communication networks,
            managing issue tracking and defect resolution through Jira during deployment phases.
          }
          \resumeItem{
            Supported integration, validation, and troubleshooting activities during commissioning
            phases in operational environments.
          }
        \resumeItemListEnd
  \resumeSubHeadingListEnd

%-----------EDUCATION-----------
\section{Education}
  \resumeSubHeadingListStart
    \resumeSubheading
      {M.Sc. in Automation Engineering}{}
      {University of Naples Federico II $|$ 110/110 cum laude}{Naples, IT}
    \resumeSubheading
      {B.Sc. in Automation Engineering}{}
      {University of Naples Federico II $|$ 107/110}{Naples, IT}
    \resumeSubheading
      {Full Stack Development \& AI Agents Specialization}{In Progress}
      {Start2Impact University (Online)}{}
  \resumeSubHeadingListEnd

%-----------TECHNICAL SKILLS-----------
\section{Technical Skills}
 \begin{itemize}[leftmargin=0.15in, label={}]
    \small{\item{
     \textbf{Embedded \& Real-Time Systems}{: Embedded C/C++, real-time systems design, memory management \& performance optimization, embedded Linux environment, debugging \& troubleshooting in field environments} \\
     \textbf{Communication \& Distributed Systems}{: Ethernet-based onboard networks, TRDP (Train Real-time Data Protocol), IPTCom, MQTT, TCP/IP networking, multi-vehicle communication logic, communication gateway design} \\
     \textbf{System Engineering \& Validation}{: System architecture definition, Integration/Verification \& Validation (IV\&V), commissioning support, requirements analysis \& traceability (DOORS), issue tracking (Jira), safety-critical environment exposure (railway domain)} \\
     \textbf{Programming Languages}{: C (Embedded), C++, Python, JavaScript, MATLAB/Simulink} \\
     \textbf{Software Development}{: Python, REST API design, PostgreSQL, Docker \& Docker Compose, Linux scripting} \\
     \textbf{Development Tools}{: CMake, Git, SVN, GitHub, Wireshark, CI/CD fundamentals, MATLAB/Simulink}
    }}
 \end{itemize}
 \textbf{Languages}: \textit{Italian} (native), \textit{English} (professional working proficiency, B2) \\
 \textbf{Availability}: Open to relocation and international assignments \\
 \textbf{Interests}: CrossFit, swing dancing, technology, graphic design, travel

%-----------PUBLICATIONS-----------
\section{Publications}
  \resumeSubHeadingPublicationsListStart
    \resumePaperTitle
      {Model-based vehicular prognostics framework using Big Data architecture}
      {Computers in Industry, Elsevier}
      {A. Petrillo, A. Picariello, S. Santini, \textbf{B. Scarciello}, G. Sperl\`i}
      {2019}
      {\fancyurl{https://doi.org/10.1016/j.compind.2019.103177}{doi.org/10.1016/j.compind.2019.103177}}
    \resumeItemListStart
      \resumeItem{
        Proposed \textbf{ATOMIC} (Automotive predicTOr Maintenance In Cloud), a cloud-based Big Data
        infrastructure for Integrated Vehicle Health Management (IVHM) that combines nonlinear
        discrete-time models of nominal and faulty component behaviour to estimate End-Of-Life (EOL)
        and Remaining Useful Life (RUL) indicators on-line.
      }
      \resumeItem{
        Validated the prognostic method on the Delphi DFG1596 fuel pump and benchmarked the
        architecture's deployment configurations in terms of scalability and cost.
      }
    \resumeItemListEnd
  \resumeSubHeadingListEnd

%-------------------------------------------
\end{document}
